\subsection{交换收敛级数}

设 \((x_{n})\) 是 \(G\) 中的可和序列，并设 \(s = \sum_{n\in \mathbb{N}}x_n\) 是它的和。那么对每个邻域 \(V\)（\(0\) 的邻域），存在 \(J_0\in \mathfrak{F}(N)\)，使得 \(s_J\in s + V\) 只要 \(J\in \mathfrak{F}(N)\) 且 \(J_0\subset J\)。设 \(m\) 是 \(J_0\) 中最大的整数；则当 \(n\geqslant m\) 时有 \(s_n\in s + V\)，因而级数 \((x_{n})\) 收敛且其和为 \(s\)。但反之不成立：

收敛级数的项所成的序列完全可能不是可和的（见第四章，§ 7）。

此外，收敛级数的定义本质上涉及 N 的序结构。若级数 \((x_{n})\) 收敛，且 \(\sigma\) 是 N 的一个置换，则级数 \((x_{\sigma(n)})\) 未必收敛（参见第四章，§ 7，习题 15）。

定义 2. 由序列 \((x_{n})\) 所定义的级数称为交换收敛的，如果对每个置换 \(\sigma\)（\(\mathbf{N}\) 的置换），由序列 \((x_{\sigma(n)})\) 所定义的级数都收敛。

命题 9. 由序列 \((x_{n})\) 所定义的级数交换收敛当且仅当序列 \((x_{n})\) 可和；且此时对每个置换 \(\sigma\)（\(\mathbf{N}\) 的置换），我们有

\[
\sum_ {n = 0} ^ {\infty} x _ {\sigma} = \sum_ {n \in \mathbb {N}} x _ {n}.
\]

若序列可和，则级数显然交换收敛。为证反向的蕴含关系，我们用反证法，设级数 \((x_{n})\) 交换收敛而序列 \((x_{n})\) 不可和。于是滤子 \(\Phi\) 在映射 \(\mathbf{H} \to s_{\mathbf{H}}\) 下的像不能是 \(\mathbf{G}\) 中的 Cauchy 滤子基，否则这个滤子基将收敛，因为按假设它有一个聚点（第二章，§ 3，第 2 节，命题 5，推论 2）。因此存在一个邻域 \(\mathbf{V}\)（\(o\) 在 \(\mathbf{G}\) 中的一个邻域），使得对每个有限子集 \(\mathbf{J}\)（\(\mathbf{N}\) 的有限子集），存在有限子集 \(\mathbf{H}\)（\(\mathbf{N}\) 的有限子集），它不与 \(\mathbf{J}\) 相交，且使得 \(\sum_{n \in \mathbf{H}} x_{n} \notin V\)。于是，用归纳法，我们可以定义

把 N 分拆为有限子集 \(H_{k}\)（ \(k \in N\)），使得对无穷多个指标 k 有 \(\sum_{n \in H_{k}} x_{n} \notin V\)。显然存在 N 的一个置换 \(\sigma\)，使得对每个 k，使 \(\sigma(n) \in H_{k}\) 的那些 n 的值是相继的。若 \(\sigma\) 是这样一个置换，则通项为 \(x_{\sigma(n)}\) 的级数不可能收敛，于是得到矛盾。

若群 G 写作乘法，由 G 的点的序列 \((x_{n})\) 所定义的无穷乘积（或通项为 \(x_{n}\) 的无穷乘积，甚或在无混淆之虞时径称乘积 \(x_{n}\)）定义为序列 \((x_{n})\) 与部分积序列 \(p_{n}=\prod_{k=0}^{n}x_{k}\) 所构成的偶。若序列 \((p_{n})\) 收敛，则称该无穷乘积收敛，这个序列的极限记作 \(\prod_{n=0}^{\infty}x_{n}\)（按记号之便，或记 \(\prod_{n=0}^{\infty}x_{n}\)）。我们把将已建立的级数性质改写为乘法记号的任务留给读者。

